\newmdenv[topline=false,leftline=false,bottomline=false]{rightborder}

\chapter{Fonctionnalités (à revoir)}

\section*{Composants (à revoir)}

Cette section liste à quoi correspond les différents icônes pouvant être rencontrés dans Slyum ainsi que le chapitre qui leur est associé.

\subsection*{Composants de Slyum}


% Projets
\begin{minipage}[c]{.2\textwidth}
	\begin{rightborder}\begin{tabular}{cc}
		\includegraphics{images/icon/open.png} &
		\includegraphics{images/icon/new.png} \\
		\includegraphics{images/icon/save.png} &
		\includegraphics{images/icon/print.png} \\
		\includegraphics{images/icon/pdf-16.png} &
		... \\
	\end{tabular}\end{rightborder}
\end{minipage}
\begin{minipage}[c]{.1\textwidth}\end{minipage}
\begin{minipage}[c]{.6\textwidth}
	\textbf{Projet} \\
	Représente les outils pour gérer les projets, les exportations ou encore les impressions.
\end{minipage}
\begin{minipage}[c]{.1\textwidth}
	\hfill
	p.{\LARGE \pageref{sec:projet}}
\end{minipage}

% Vues
\begin{minipage}[c]{.2\textwidth}
    \begin{rightborder}\begin{tabular}{cc}
		\includegraphics{images/icon/element-view.png} &
		\includegraphics{images/icon/element-view-open.png} \\
		\includegraphics{images/icon/element-view-add.png} &
		\includegraphics{images/icon/element-view-delete.png} \\
    \end{tabular}\end{rightborder}
\end{minipage}
\begin{minipage}[c]{.1\textwidth}\end{minipage}
\begin{minipage}[c]{.6\textwidth}
    \textbf{Vues} \\
    Représente une vue du projet. Permet de créer, ouvrir et supprimer des vues.
\end{minipage}
\begin{minipage}[c]{.1\textwidth}
    \hfill
    p.{\LARGE \pageref{sec:vues}}
\end{minipage}

\subsection*{Composants UML}

% Entités
\begin{minipage}[c]{.2\textwidth}
    \begin{rightborder}\begin{tabular}{cc}
        \includegraphics{images/icon/class.png} &
        \includegraphics{images/icon/interface.png} \\
        \includegraphics{images/icon/enum.png} &
        \includegraphics{images/icon/classAssoc.png} \\
    \end{tabular}\end{rightborder}
\end{minipage}
\begin{minipage}[c]{.1\textwidth}\end{minipage}
\begin{minipage}[c]{.6\textwidth}
    \textbf{Entités} \\
    Représente les différentes entités UML telles que les classes ou les interfaces.
\end{minipage}
\begin{minipage}[c]{.1\textwidth}
    \hfill
    p.{\LARGE \pageref{sec:entites}}
\end{minipage}

% Association
\begin{minipage}[c]{.2\textwidth}
    \begin{rightborder}\begin{tabular}{cc}
        \includegraphics{images/icon/generalize.png} &
        \includegraphics{images/icon/dependency.png} \\
        \includegraphics{images/icon/association.png} &
        \includegraphics{images/icon/aggregation.png} \\
        \includegraphics{images/icon/composition.png} &
        \includegraphics{images/icon/innerClass.png} \\
        \multicolumn{2}{c}{\includegraphics[width=13px]{images/icon/multi.png}} \\
    \end{tabular}\end{rightborder}
\end{minipage}
\begin{minipage}[c]{.1\textwidth}\end{minipage}
\begin{minipage}[c]{.6\textwidth}
    \textbf{Relations} \\
    Représente les différentes relations entre entités UML telle que les associations, les aggrégations ou encore les héritages.
\end{minipage}
\begin{minipage}[c]{.1\textwidth}
    \hfill
    p.{\LARGE \pageref{sec:relations}}
\end{minipage}

% Notes
\begin{minipage}[c]{.2\textwidth}
	\begin{rightborder}\begin{tabular}{cc}
		\includegraphics{images/icon/note.png} &
		\includegraphics{images/icon/linkNote.png} \\
		\multicolumn{2}{c}{\includegraphics{images/icon/multiNote.png}} \\
	\end{tabular}\end{rightborder}
\end{minipage}
\begin{minipage}[c]{.1\textwidth}\end{minipage}
\begin{minipage}[c]{.6\textwidth}
	\textbf{Notes} \\
	Représentation des notes. Permet de créer des notes sur le diagramme de classe.
\end{minipage}
\begin{minipage}[c]{.1\textwidth}
	\hfill
	p.{\LARGE \pageref{sec:notes}}
\end{minipage}

% Outils de personnalisation
\begin{minipage}[c]{.2\textwidth}
	\begin{rightborder}\begin{tabular}{cc}
		\includegraphics{images/icon/pointer-grip.png} &
		\includegraphics{images/icon/alignLeft.png} \\
		\includegraphics{images/icon/delete.png} &
		\includegraphics{images/icon/color.png} \\
		\includegraphics{images/icon/top.png} &
		... \\
	\end{tabular}\end{rightborder}
\end{minipage}
\begin{minipage}[c]{.1\textwidth}\end{minipage}
\begin{minipage}[c]{.6\textwidth}
	\textbf{Outils de personnalisation} \\
	Représente les outils permettant de personnaliser et modifier les différents composants du diagramme.
\end{minipage}
\begin{minipage}[c]{.1\textwidth}
	\hfill
	p.{\LARGE \pageref{ch:outils}}
\end{minipage}